% Section 3: the cost side. Two housing figures in this section look
% contradictory unless the text says which measure it is using and where each
% measure comes from, so the reconciliation is explicit and every rent source is
% named with its universe and vintage at first mention.
\section{Austin is an average-cost city with expensive housing}
\label{sec:costs}

What Austin musicians earn matters only against what living in Austin costs.
Austin is often described in local reporting and in national coverage as an
expensive city across the board, from groceries to rent to services, but the
federal price record narrows that picture: Austin runs close to the national
average on most goods and services and above the national average only in
housing.\footnote{Framings that treat Austin as broadly expensive appear in
national coverage such as U.S.\ News \& World Report's cost-of-living
rankings and in Austin American-Statesman and Axios Austin coverage of
city-affordability trends; this study does not enumerate that literature
here, and the BEA Regional Price Parities cited in the next sentence are
the federal counterpoint.} The Bureau of
Economic Analysis publishes an annual price comparison across metro areas
called the Regional Price Parities (RPP), which measures the price of a common
basket of goods and services in each metro against a national average set to
100 and reports it separately for four component baskets, one of which is
housing services.\footnote{The BEA program is the \textit{Regional Price
Parities and Implicit Regional Price Deflators by MSA}, currently released in
its 2024 vintage covering years 2008 through 2024, and available from the
Bureau of Economic Analysis at
\url{https://apps.bea.gov/regional/zip/MARPP.zip}. The universe is the
population of each metropolitan statistical area, and the basket weights come
from the Consumer Expenditure Survey. Each metro's all-items index is a
weighted average of the four component indexes, so the housing component and
the all-items index sit on the same national base of 100 and their difference
in the same metro is in index points rather than percent. BEA publishes no
sampling error for these estimates because the underlying price series are
themselves published indexes rather than sample statistics.} On the 2024 BEA
figures, Austin's all-items index is \NColRppAllitemsAustin\ against the
national average of 100, marginally below the national average across all
goods and services. Austin's housing services index is
\NColRppHousingAustin, so housing costs in Austin run
\NColRppGapAustin\ higher than everything else does in the same city. That
place gap is the widest of the five metros compared here, ahead of
Dallas \NColRppGapDallas, Nashville \NColRppGapNashville, Houston
\NColRppGapHouston, and San Antonio essentially none at
\NColRppGapSanantonio.\footnote{The gap reported for each metro is
\(g_m = \mathrm{RPP}^{\mathrm{hous}}_m - \mathrm{RPP}^{\mathrm{all}}_m\), the
difference in index points between that metro's housing parity and its own
all-items parity, so it measures how much more expensive housing is than
everything else in the same place rather than how expensive that place is
against another one. Because both terms are indexed to the same national base
of 100, the difference is in index points and is not a percentage. Austin's
housing parity peaked at \NColRppAustinHousingPeak, a peak year the two
independent housing series in Figure~\ref{fig:metrorent} corroborate.}

The national price record shows the same relative movement over time, in a
different measure. The Bureau of Labor Statistics tracks the price of a fixed
consumer basket each month in its Consumer Price Index for All Urban Consumers
(CPI-U), for a universe of the urban population of the United States, and
reports the all-items index and its components on a common
base.\footnote{The series are the CPI-U All items (series identifier
CUUR0000SA0) and the CPI-U Housing component, U.S.\ city average, monthly, and
this study rebases each to its 2005 annual average and takes the 2025 annual
average of each rebased series. Program documentation is available from BLS at
\url{https://www.bls.gov/cpi/}. BLS publishes no Austin-area CPI at all, so
every real-dollar figure in this report deflates by the national U.S.\ city
average CPI-U, and given the RPP evidence above that Austin housing runs above
the national norm, this national deflator most likely understates the true
Austin price rise. The 2025 annual average of the housing component uses an
interpolated October, because the federal funding lapse suspended BLS price
collection that month.} Rebased so that each series equals 100 in 2005, the
housing component reaches \NColCpiHousingIndexTwoZeroTwoFive\ by 2025 against
\NColCpiAllitemsIndexTwoZeroTwoFive\ for all items, so national housing prices
have risen \NColCpiHousingGapTwoZeroTwoFive\ faster than the broad basket since
2005.\footnote{Rebasing on 2005 puts \(I^{s}_{t} = 100 \times C^{s}_{t} /
C^{s}_{2005}\), where \(C^{s}_{t}\) is the annual average of the CPI-U for
series \(s\), either all items or the housing component, in year \(t\).
Putting both series on the same base year is what makes them comparable, and
the gap is the difference of the two rebased levels in 2025,
\(I^{\mathrm{hous}}_{2025} - I^{\mathrm{all}}_{2025}\), again in index points
rather than percent. This national CPI comparison and the BEA metro parity
above answer different questions and are not comparable in level: the
CPI-Housing figure describes the national price rise for housing since 2005,
and the RPP-Housing figure above describes where the Austin metro sits against
the current national average.}

For Austin musicians the cost problem concentrates in housing, the one BEA
category that prices well above the national norm in Austin. The rest of this
section follows housing: what market rents have done since 2015, how many
hours of work a month of rent costs at Austin wages, and why those two figures
appear to disagree.

\subsection{Rents rose sharply, then fell further than in any comparison metro}

\begin{figure}[htbp]
\centering
\includegraphics[width=\textwidth]{fig20_metro_rent.png}
\caption{Real market rent indexed to each metro's own January 2015 level, 2025
dollars, through June 2026. Source: Zillow Observed Rent Index (ZORI),
seasonally adjusted, deflated with the national Consumer Price Index for All
Urban Consumers. How to read it: each line shows how far that metro's real
rent has moved from its own 2015 starting point, so the lines compare
trajectories rather than rent levels. Note on timing: Austin's real peak is
December 2021, about a year before the widely quoted nominal peak, because
deflating during the sharp inflation of 2021 to 2023 shifts the turning point
earlier. Both readings are correct and answer different questions. Caveat: the
small step every line takes at each January is the annual deflator changing,
not a change in rents, because a single year-average price index is applied to
twelve months of rent data.}
\label{fig:metrorent}
\end{figure}

Zillow's Observed Rent Index estimates the typical asking rent in each metro
each month by weighting an internal repeat-rent index against local rental
market shares, and it is the closest thing in the public record to a monthly
Austin rent series, because BLS publishes no Austin-area consumer price
index.\footnote{Zillow Group publishes the seasonally adjusted metro-level
ZORI series at \url{https://www.zillow.com/research/data/}, together with a
methodology document. The universe is rental listings on the Zillow platform,
weighted to the U.S.\ Postal Service's count of the rental housing stock in
each geography, so ZORI captures the rent asked on newly available units
rather than the rent paid on all standing leases. It is not a probability
sample of dwellings, and Zillow does not attach a sampling error to it. It
runs monthly at the metro level and picks up cyclical turns faster than the
Census Bureau's American Community Survey median gross rent, which pools five
years of household responses and moves slowly by construction.} On that
series, Austin's real market rent peaked at \NColZoriAustinPeak\ and has
fallen \NColZoriAustinPeaktolatest\ since, the steepest decline among the five
metros in Figure~\ref{fig:metrorent}, with San Antonio the nearest at
\NColZoriSanantonioPeaktolatest\ and Dallas
\NColZoriDallasPeaktolatest\ and Houston
\NColZoriHoustonPeaktolatest.\footnote{The plotted series deflates each metro
first and rebases second. The real level is
\(R^{\mathrm{real}}_{m,t} = R_{m,t} \times d_{y(t)}\), where \(R_{m,t}\) is
the seasonally adjusted ZORI value for metro \(m\) in month \(t\), \(y(t)\)
is the calendar year containing that month, and \(d_y = C_{2025}/C_y\) is
the ratio of national CPI-U annual averages that carries year \(y\) dollars
to 2025 dollars. The plotted line is then
\(I_{m,t} = 100 \times R^{\mathrm{real}}_{m,t} \big/
   R^{\mathrm{real}}_{m,\text{Jan }2015}\),
each metro against its own base month, which removes the level and leaves the
trajectory, so metros at different rent levels can share one axis. The
peak-to-latest change is
\(100\,\bigl(R^{\mathrm{real}}_{m,\text{latest}} \big/
   R^{\mathrm{real}}_{m,\text{peak}} - 1\bigr)\)
on the deflated series. Deflating with an annual CPI-U during the sharp
inflation of 2021 to 2023 is what moves the real peak about a year ahead of
the nominal one: the step in \(d_y\) between 2021 and 2022 outweighs the
further nominal rent increase in that year.} Home values fell further still.
Zillow's Home Value Index (ZHVI), a repeat-sales estimate of the typical
market value of a single-family home in the metro, records an Austin peak of
\NColZhviAustinPeak\ and a real decline of
\NColZhviAustinPeaktolatest\ from it.\footnote{Zillow Group publishes the
metro-level ZHVI (all homes, seasonally adjusted) at the same research portal,
\url{https://www.zillow.com/research/data/}. The universe is the middle tier
of homes valued between the 35th and 65th percentiles in each geography,
estimated monthly from Zillow's proprietary automated valuation model and
transaction records. ZHVI is a market-value estimate rather than a
transaction-price series, so it can move faster than recorded sales prices
where turnover thins, and no sampling error attaches to it.} A report that
showed only the run-up and stopped in 2022 would be describing a market that
no longer exists.

The federal housing-assistance benchmark moved differently. The Department
of Housing and Urban Development publishes an annual Fair Market Rent (FMR)
for each metropolitan area, one figure per bedroom count, and sets it near
the 40th percentile of standard-quality rents in the area, using a rolling
five-year window of American Community Survey responses updated with more
recent private data.\footnote{The all-years FMR history is available from
HUD at \url{https://www.huduser.gov/portal/datasets/FMR/FMR_All_1983_2026.csv}.
The FMR is the reference point for what the federal Housing Choice Voucher
program, and other HUD rental-assistance programs, will pay in each metro
area. It is a benchmark, not an observed rent: HUD constructs it from ACS
five-year survey data on gross rent for standard-quality units and then
trends it forward with more recent private rent indexes. The universe is
recent movers within the metro, so FMR turns later than a market series like
ZORI. Two breaks in the source sit inside this study's window and make the
peak date and the peak-to-latest change approximate: the nominal one-bedroom
Austin rent jumps 18.0 percent between FY2023 and FY2024, the same year the
file's percentile-basis field goes blank, and the FY2026 value is published
for a redefined metro area, Austin-Round Rock-San Marcos rather than
Austin-Round Rock.} HUD's Austin one-bedroom FMR peaked later, at
\NColHudBrOneAustinPeak, and has come down only
\NColHudBrOneAustinPeaktolatest\ since.\footnote{Both housing series in this
subsection use the same deflator and dating rule as the market rent above,
\(F^{\mathrm{real}}_y = F_y \times d_y\) followed by
\(100\,(F^{\mathrm{real}}_{\mathrm{latest}}/F^{\mathrm{real}}_{\mathrm{peak}}
   - 1)\), so both are in constant 2025 dollars. The hours-of-work ratio in
the next subsection uses the same HUD Fair Market Rent in nominal dollars
instead, because there the rent is divided by a same-year wage and deflating
both sides would cancel.} Austin's assistance benchmark turned about two and
a half years after the market series did. That lag matters for any program
priced off the benchmark rather than off the market, and it is the larger
part of the tension between the next two figures. The rest of that tension is
that the payroll musician wage the next figure divides by fell back from an
unusually high 2022 reading in the years HUD's FMR was still rising.

\subsection{The payroll wage understates what rent takes from musicians}

\begin{figure}[htbp]
\centering
\includegraphics[width=\textwidth]{fig19_rent_vs_wages.png}
\caption{Hours of work at the median hourly wage needed to cover one month of
Austin one-bedroom Fair Market Rent. Sources: Bureau of Labor Statistics
Occupational Employment and Wage Statistics for the Austin metro; HUD Fair
Market Rents; and the City of Austin's purchased earnings series, dataset
\texttt{2qxc-8cme}. Wages and rents are matched within year in nominal terms,
because the ratio of two same-year dollar figures needs no inflation
adjustment. How to read it: a higher line means rent costs more hours of work.
Critical caveat: the federal wage series counts payroll jobs only and excludes
the self-employed, who are \NPumsAustinMusSelfemp\
(\NPumsAustinMusSelfempMoe) of employed Austin-metro musicians aged 18 to 64
over 2020 to 2024, and the Austin musician estimate rests on between about
100 and 700 payroll jobs depending on the year, and
\NEmpPayrollMusiciansAustinLast\ in 2025, which makes it volatile from year to
year. The city's own purchased earnings series, which appears to include
self-employment income, exists for only two years and is not plotted; in 2017
it put musician pay at \NColAnchorCityWageVsOewsTwoZeroOneSeven\ of the
federal figure for the same occupation.}
\label{fig:rentwages}
\end{figure}

The federal payroll wage this figure divides into rent is the Austin-metro
median hourly wage in the BLS Occupational Employment and Wage Statistics
program (OEWS), a semi-annual establishment survey that BLS and the state
workforce agencies run jointly.\footnote{BLS publishes the OEWS Austin
Metropolitan Statistical Area estimates in its annual tables at
\url{https://www.bls.gov/oes/tables.htm}, the most recent vintage being
May~2025. Establishments in the sample report the wage of every worker on
their payroll but report nothing for the self-employed, who fall outside the
sampling frame, so the Austin musician wage in this figure describes only the
\NEmpPayrollMusiciansAustinLast\ musician jobs the 2025 OEWS Austin sample
found on a payroll, and Section~\ref{sec:measurement} develops that
measurement gap in full. BLS does not publish an \emph{annual} musician wage
for any city or year because musician work is too irregular to annualise from
an hourly rate, so the ratio here uses the hourly median.} On that payroll
wage, a month of one-bedroom rent costs an Austin musician
\NColAnchorHoursMusicianLatest\ of work, against
\NColAnchorHoursAlloccLatest\ for the median Austin worker in 2025
(Figure~\ref{fig:rentwages}).\footnote{The ratio is
\(h_{o,y} = F_y / w_{o,y}\), where \(F_y\) is the Austin one-bedroom HUD Fair
Market Rent for fiscal year \(y\) in dollars per month, \(w_{o,y}\) is the
Austin-metro OEWS median hourly wage for occupation group \(o\) in the same
year in dollars per hour, and \(h_{o,y}\) is therefore in hours per month.
This ratio needs no inflation adjustment, and both sides are left in nominal
dollars on purpose: multiplying numerator and denominator by the same factor
\(d_y\) leaves \((d_y F_y)/(d_y w_{o,y}) = F_y / w_{o,y}\), so deflating would
change nothing except to invite a mismatch of years. The two sides are
matched within year for the same reason. BLS did not publish the musician
wage for 2016 and the figure draws a straight segment across that year, so
the 2015-to-2017 stretch of the orange line should not be read as an observed
path. This is a ratio of two published medians, not a sample estimate, so
it has no standard error; its year-to-year movement still reflects sampling
noise in the thin musician wage estimate underneath it.} By that measure
Austin musicians look better placed than the typical Austin worker, because
the OEWS payroll wage covers only the \NEmpPayrollMusiciansAustinLast\
musician jobs the 2025 OEWS Austin sample recorded, the measurement gap
Section~\ref{sec:measurement} describes, and those payroll jobs pay above the
Austin all-occupation median. That ordering is a real feature of the payroll
data and a poor description of the occupation, and on so thin a base the
shape of the orange line across the full period is firmer evidence than the
2025 point on its own.

The picture changes when the earnings measure includes the self-employed. The
Census Bureau's American Community Survey Public Use Microdata Sample (ACS
PUMS) samples households rather than employers, so it captures earnings for
the self-employed alongside payroll earners, at the cost of thinner samples in
any single year.\footnote{The ACS is an ongoing annual household survey the
Census Bureau administers to approximately 3.5 million addresses each year;
the five-year PUMS pools respondents across 2020 through 2024 to produce
metro-level estimates on smaller occupations, at
\url{https://www2.census.gov/programs-surveys/acs/data/pums/2024/5-Year/}. The
Austin-metro musician cell used here is Musicians and Singers (occupation
code OCCP 2752), employed 18 to 64 with positive earnings, weighted by PWGTP.
Because the five-year file is one pooled cross-section rather than five
annual observations, this share is a point estimate for the pooled period
rather than a series that can be plotted year by year.} Measured against ACS
median earnings from Section~\ref{sec:earnings}, a year of one-bedroom HUD
Fair Market Rent would equal \NColPumsRentShareMatched\ of the median annual
earnings of employed Austin-metro musicians aged 18 to 64 over 2020 to
2024.\footnote{The share is
\(s = 100 \times 12\bar{F} \big/ \tilde{y}\), where \(\bar{F}\) is the
average Austin one-bedroom HUD Fair Market Rent across fiscal years 2020 to
2024 in 2025 dollars, the same five years the microdata pool, and
\(\tilde{y}\) is the weighted median personal earnings from work of Austin
musicians in 2025 dollars on the positive-earnings base. Averaging the five
fiscal years rather than picking one is what makes the numerator and the
denominator cover the same period. The numerator is a housing cost for a
dwelling and the denominator is one person's earnings, so \(s\) is
affordability arithmetic rather than an observed housing outlay, and it
carries the margin of error of the earnings median in its denominator.
Measured against the FY2026 rent, which replaces \(\bar{F}\) with that single
year, it is \NColPumsRentShareLatest. Both are far above the 30 percent share
the federal government treats as the ceiling for a rental burden a household
can carry without cost-burden status, which is a statement about arithmetic
rather than a claim about any individual's housing.} The numerator here is a
benchmark rent rather than an observed one. Separate corroboration comes from
Section~\ref{sec:earnings}, which reports that across Texas, renting
households with a musician in them pay a higher median gross rent than
renting worker households generally, while carrying about the same
rent-to-income burden. On that evidence, Texas musician households are not
renting more cheaply than other working households.

\subsection{Real market rents fell while the nominal assistance benchmark kept climbing through 2025}

Figure~\ref{fig:metrorent} shows real rents falling. Figure~\ref{fig:rentwages}
shows hours of work per month of rent rising between 2023 and 2025. Both are
correct, and a reader who noticed the tension without an explanation would be
right to distrust what follows. The two figures rest on different rents and
different denominators, as the table below sets out:

\begin{center}
\footnotesize
\begin{tabular}{@{}B{2.9cm} P{6.3cm} P{6.6cm}@{}}
\toprule
& \textbf{Figure~\ref{fig:metrorent}} & \textbf{Figure~\ref{fig:rentwages}} \\
\midrule
Rent measure & Zillow Observed Rent Index, a market listings series deflated
to constant 2025 dollars and indexed to each metro's own January 2015 level &
HUD Fair Market Rent for a one-bedroom in Austin, nominal dollars, the annual
benchmark for federal housing-assistance payments \\
\rowrule
Wage denominator & None: the line plots rent alone & The BLS OEWS median
hourly wage across the \NEmpPayrollMusiciansAustinLast\ Austin musician jobs
recorded in 2025 \\
\rowrule
Direction, 2023 to 2025 & Down: Zillow market rents fell in real terms in
Austin more than in any other metro compared & Up: HUD's benchmark kept
climbing while the OEWS musician wage fell back from an unusually high 2022
reading \\
\bottomrule
\end{tabular}
\end{center}

\vspace{3pt}
Where this study needs a single number for what renting costs in Austin, it
takes the Zillow market series and says so. The hours-of-rent figure also
rests on an OEWS payroll wage that describes fewer musician jobs than it once
did, and Section~\ref{sec:payroll} follows that count across twenty years.
